%Una explicaci\'on breve del contenido del informe, es decir, detalla: Prop\'osito, Estructura del Documento, Descripci\'on (muy breve) del Problema y Motivaci\'on.

%En este informe se realizará una presentación de los trabajos más importantes respecto al problema de Truck Drone Parcel Delivery Problem (TDPDP), el cual consiste en encontrar una ruta para un camión que haga entregas a ciertas ubicaciones en el menor tiempo posible, teniendo el apoyo de un dron que permite hacer cargas en paralelo o con menor tiempo respecto al camión.

%La motivación de este trabajo nace de análisis como \cite{moshrefjavadi2021applications} donde se explica el gran auge que ha tenido la utilización de drones en el último tiempo, especialmente con motivos de entrega de productos, dado su alto valor económico para las empresas.

%A continuación, se presentarán los principales avances desde su concepción inicial en el trabajo de Murray et al. (2014) \cite{Murray2015Flying}, pasando también por variantes definidas en otros trabajos, y comparando de qué manera se resolvían estos problemas en computación.

En este informe se realizará una presentación de la implementación de un algoritmo evolutivo para el problema del \textit{Truck Drone Parcel Delivery Problem} (TDPDP). El TDPDP consiste en encontrar una ruta para un camión que realice entregas a ciertas ubicaciones en el menor tiempo posible, teniendo el apoyo de un dron que permite realizar entregas en paralelo junto al camión.

La motivación de este trabajo nace de análisis como \cite{moshrefjavadi2021applications}, donde se explica el gran auge que ha tenido la utilización de drones en el último tiempo, especialmente con motivos de entrega de productos, dado su alto valor económico para las empresas.

Primero, se hará una definición del problema a tratar, explicando las características formales del problema, además de tratar su complejidad y variantes. A continuación, se presentará un análisis del estado del arte del problema, comparando con los trabajos principales tales como el trabajo seminal de \textit{Murray} y \textit{Chu} \cite{Murray2015Flying}. Tras esto, se presenta una definición formal del modelo matemático y la representación que se utilizará para el algoritmo, especificando variables, restricciones, variantes y función objetivo.

Con el modelo y la representación definidos, se procederá a la descripción de la implementación del algoritmo evolutivo, explicando cómo opera bajo el enfoque \textit{Truck-First, Drone-Second}. Con esto, se explicarán los experimentos realizados con el algoritmo y los resultados que estos entregaron. Finalmente, se presentará una conclusión sobre el estado del problema, sobre los resultados obtenidos de la implementación evolutiva con enfoque \textit{Truck-First, Drone-Second}, y sobre trabajos e implementaciones para el futuro.
